% ============================================================================
% Appendices: optional mechanism and supporting diagnostics
% ============================================================================
\appendices

\section{Optional Bounded Preference Coupling}
\label{app:coupling}

Accepted regional anchors can additionally be translated into a bounded
adjustment of the acquisition mean.
This posterior-mean coupling is a diagnostic variant that is separate from the
acquisition-restart regional guidance used by the principal LLMBO-MO
configuration; it is therefore analyzed independently here.

Let $\bG_t=\{\mathbf g_{j,t}\}_{j=1}^{M}$ denote accepted anchors.
Each anchor is scored using its EI value, posterior uncertainty, and novelty:
\begin{equation}
\begin{aligned}
n_{j,t}
&=\min_{\mathbf x\in\bD_t}
\|\widehat{\mathbf g}_{j,t}-\widehat{\mathbf x}\|_2,\\
s_{j,t}
&=\eta_{\mathrm{EI}}z\!\left(
\log\max\{\mathrm{EI}_{j,t},\varepsilon_{\mathrm{EI}}\}\right)
+\eta_\sigma z(\sigma_{j,t})+\eta_n z(n_{j,t}),\\
v_{j,t}
&=\frac{\exp(s_{j,t}/T_v)}
{\sum_{\ell=1}^{M}\exp(s_{\ell,t}/T_v)}.
\end{aligned}
\label{eq:anchor_weights}
\end{equation}
If anchor scoring is unavailable numerically, uniform weights
$v_{j,t}=1/M$ are used.

Let $\boldsymbol\Sigma_{GG,t}$ be the posterior covariance among anchors and
$\boldsymbol\Sigma_{\theta G,t}$ the candidate-to-anchor covariance.
The posterior correlation between a candidate and the weighted anchor portfolio
is
\begin{equation}
\begin{aligned}
V_{G,t}
&=\mathbf v_t^\top
(\boldsymbol\Sigma_{GG,t}+\varepsilon_K\mathbf I)\mathbf v_t,\\
\rho_t(\btheta)
&=\frac{\boldsymbol\Sigma_{\theta G,t}\mathbf v_t}
{\sigma_t(\btheta)\sqrt{\max\{V_{G,t},\varepsilon_v\}}}.
\end{aligned}
\label{eq:covariance_projection}
\end{equation}
Only positive correlations are used to attract EI toward preferred regions.

The region-quality gate combines the concentration of anchor weights and the
dispersion of anchor scores:
\begin{equation}
\begin{aligned}
H_t&=-\sum_{j=1}^{M}v_{j,t}\log(v_{j,t}+\varepsilon_H),\\
\bar H_t&=\frac{H_t}{\max\{\log M,\varepsilon_H\}},\\
q_t&=\operatorname{clip}\!\left[
\eta_H(1-\bar H_t)
+(1-\eta_H)\tanh\!\bigl(\operatorname{std}(\mathbf s_t)\bigr),
0,1\right].
\end{aligned}
\label{eq:quality_gate}
\end{equation}
The LLM-reported confidence
$c_t=\operatorname{clip}(h_t^{\mathrm{LLM}},0,1)$ and a bounded trust factor
$\tau_t\in[0,1]$ define $r_t=c_t\tau_tq_t$.
The attenuation $a_t=\operatorname{clip}(1-(t-1)/T_{\mathrm L},0,1)$ reduces
the possible influence to zero by the end of the guidance window.
The bounded adjustment is
\begin{equation}
\begin{aligned}
\delta_t(\btheta)
&=\operatorname{clip}\!\left(
a_t\lambda_{\max}r_t[\rho_t(\btheta)]_+,
0,\delta_{\max}\right),\\
\mu_t^{\mathrm L}(\btheta)
&=\mu_t(\btheta)-\delta_t(\btheta),\\
(\sigma_t^{\mathrm L})^2(\btheta)&=\sigma_t^2(\btheta).
\end{aligned}
\label{eq:bounded_coupling}
\end{equation}
Because \eqref{eq:tchebycheff} is minimized, decreasing the posterior mean near
positively correlated anchors increases their attractiveness to EI.
Clipping limits the adjustment in standardized scalar-target units, preserves
the posterior variance, and creates no pseudo-observation.
The principal results use regional information only as acquisition-search
restarts and therefore do not rely on \eqref{eq:bounded_coupling}.

\section{Supporting Numerical Diagnostics}
\label{app:diagnostics}

\subsection{Pareto-Archive Growth}

The archive-growth diagnostic tracks the number of feasible
Pareto-nondominated protocols retained during optimization.
It describes archive richness and is not used as a substitute for HV or as a
primary sample-efficiency metric.
Here, an ``optimal charging protocol'' means a feasible Pareto-nondominated
entry in the current archive.

\begin{figure}[H]
\centering
\includegraphics[width=0.99\columnwidth]{optimal_protocols_curve_600dpi.png}
\caption{Chen2020 Pareto-archive growth over 50 algorithmic iterations. Lines
are five-seed means and bands show across-seed dispersion.}
\label{fig:optimal_protocols}
\end{figure}

At iteration 50, LLMBO-MO maintains an average of approximately
\OptimalLLMBOFinal\ nondominated protocols, compared with
\OptimalParEGOFinal\ for ParEGO.
Its curve remains higher through most of the search, indicating a richer
retained set of trade-offs across these seeds.
Because the horizontal axis counts algorithmic iterations,
Fig.~\ref{fig:optimal_protocols} is not an equal-simulator-call sample-efficiency
comparison and is independent of the fixed-budget HV result.

\subsection{Objective-Scaling Sensitivity}

The scaling diagnostic compares min--max scaling, Z-score standardization, and
no normalization on Chen2020 over five seeds and the same 56-evaluation budget.
It tests numerical robustness rather than the LLM mechanism.

\begin{figure}[H]
\centering
\includegraphics[width=0.99\columnwidth]{objective_normalization_hv_600dpi.png}
\caption{Chen2020 objective-preprocessing sensitivity over five seeds. Lines
show mean HV and bands show one sample standard deviation.}
\label{fig:normalization_hv}
\end{figure}

\begin{table}[!t]
\caption{Objective-Preprocessing Sensitivity on Chen2020}
\label{tab:normalization_hv}
\centering
\footnotesize
\setlength{\tabcolsep}{3.8pt}
\begin{tabular}{@{}lcc@{}}
\toprule
Preprocessing & Evaluation 30 & Evaluation 56 \\
\midrule
Min--max & \NormMidMinMax & \NormMinMaxMean\,$\pm$\,\NormMinMaxStd \\
Z-score & \textbf{\NormMidZScore} & \NormZScoreMean\,$\pm$\,\NormZScoreStd \\
None & \NormMidNone & \NormNoneMean\,$\pm$\,\NormNoneStd \\
\bottomrule
\end{tabular}
\end{table}

Z-score preprocessing progresses faster in the middle of the search, reaching
an HV of \NormMidZScore\ at evaluation 30 compared with \NormMidMinMax\ for
min--max scaling.
At evaluation 56, min--max scaling has the slightly higher mean
(\NormMinMaxMean\ versus \NormZScoreMean), but the difference is small and is
not interpreted as a significant ordering.
Both scale-balancing choices have substantially lower final dispersion than
the unnormalized condition, whose final result is
\NormNoneMean\,$\pm$\,\NormNoneStd.
Thus, heterogeneous objective magnitudes should be balanced, while either
conventional scaling choice is defensible in the present setting.


\section{Seed Configuration and Reproducibility}
\label{app:reproducibility}

The main Chen2020 and Ecker2015 comparisons use the same five seeds
(\MainSeedSet) for each method.
The reported studies share a common configuration:

\begin{itemize}
\item Chen2020 main comparison: five seeds (\MainSeedSet);
56 simulator calls per seed.
\item Ecker2015 second-parameterization case: five seeds
(\MainSeedSet); 56 simulator calls per seed.
\item Initialization study: ten matched seeds; 26-call budget.
\item Component ablation: three matched seeds (8409--8411);
16-call budget.
\item Main regional guidance: $T_{\mathrm L}=16$, 64 anchors,
EI-search restarts.
\item LLM records: prompt configuration and responses logged per
seed.
\end{itemize}

Seed-level identifiers, LLM call records, and parsed responses are
retained in the experiment logs rather than in the main parameter
table so that the reported statistics can be traced to the
corresponding seed set without interrupting the main experimental
narrative.

For the main algorithm, regional information affects only
acquisition-search restarts. The posterior-mean coupling in
Appendix~\ref{app:coupling} is a separate diagnostic realization and
is not used to define the principal results.
